\documentclass[10pt,a4paper]{amsart}
\usepackage{amsmath,amssymb,amsthm,graphicx}
% Retained preamble for the Stacks Project chapter snapshot in this directory.
% The chapter's own preamble.tex is a copy of the shared Stacks Project preamble; the block
% below is sliced MECHANICALLY, line for line, from it. Nothing was added, removed or
% redefined; the wrapper class, the xr-hyper/hyperref wiring and the externaldocument list
% are omitted because the assembler supplies the document class.
%
% lines taken from preamble.tex: 3, 5, 9, 11, 12, 156, 157, 158, ... (51 lines total)
%
% package coverage restored from the chapter's own preamble (2026-10-13): xy, xy-2cell

\usepackage[all]{xy}\usepackage{verbatim}
\xyoption{2cell}
\UseAllTwocells

\newenvironment{reference}{\comment}{\endcomment}
\newenvironment{slogan}{\comment}{\endcomment}
\newenvironment{history}{\comment}{\endcomment}
\theoremstyle{plain}
\newtheorem{theorem}[subsection]{Theorem}
\newtheorem{proposition}[subsection]{Proposition}
\newtheorem{lemma}[subsection]{Lemma}
\theoremstyle{definition}
\newtheorem{definition}[subsection]{Definition}
\newtheorem{example}[subsection]{Example}
\newtheorem{exercise}[subsection]{Exercise}
\newtheorem{situation}[subsection]{Situation}
\theoremstyle{remark}
\newtheorem{remark}[subsection]{Remark}
\newtheorem{remarks}[subsection]{Remarks}
\numberwithin{equation}{subsection}
\def\lim{\mathop{\mathrm{lim}}\nolimits}
\def\colim{\mathop{\mathrm{colim}}\nolimits}
\def\Spec{\mathop{\mathrm{Spec}}}
\def\Hom{\mathop{\mathrm{Hom}}\nolimits}
\def\Ext{\mathop{\mathrm{Ext}}\nolimits}
\def\SheafHom{\mathop{\mathcal{H}\!\mathit{om}}\nolimits}
\def\SheafExt{\mathop{\mathcal{E}\!\mathit{xt}}\nolimits}
\def\Sch{\mathit{Sch}}
\def\Mor{\mathop{\mathrm{Mor}}\nolimits}
\def\Ob{\mathop{\mathrm{Ob}}\nolimits}
\def\Sh{\mathop{\mathit{Sh}}\nolimits}
\def\NL{\mathop{N\!L}\nolimits}
\def\CH{\mathop{\mathrm{CH}}\nolimits}
\def\proetale{{pro\text{-}\acute{e}tale}}
\def\etale{{\acute{e}tale}}
\def\QCoh{\mathit{QCoh}}
\def\Ker{\mathop{\mathrm{Ker}}}
\def\Im{\mathop{\mathrm{Im}}}
\def\Coker{\mathop{\mathrm{Coker}}}
\def\Coim{\mathop{\mathrm{Coim}}}
\DeclareMathSymbol{\boxtimes}{\mathbin}{AMSa}{"02}
\def\QCohstack{\mathcal{QC}\!\mathit{oh}}
\def\Cohstack{\mathcal{C}\!\mathit{oh}}
\def\Spacesstack{\mathcal{S}\!\mathit{paces}}
\def\Quotfunctor{\mathrm{Quot}}
\def\Hilbfunctor{\mathrm{Hilb}}
\def\Curvesstack{\mathcal{C}\!\mathit{urves}}
\def\Polarizedstack{\mathcal{P}\!\mathit{olarized}}
\def\Complexesstack{\mathcal{C}\!\mathit{omplexes}}
\def\Pic{\mathop{\mathrm{Pic}}\nolimits}
\def\Picardstack{\mathcal{P}\!\mathit{ic}}
\def\Picardfunctor{\mathrm{Pic}}
\def\Deformationcategory{\mathcal{D}\!\mathit{ef}}
\usepackage{cleveref}
\crefname{originallemma}{original lemma}{original lemmas}
\Crefname{originallemma}{Original lemma}{Original lemmas}
\makeatletter
\newlabel{cohomology-lemma-pushforward-restriction}{{08DF}{source}{}{sourcepointer-1}{}}
\newlabel{cohomology-lemma-pushforward-restriction@cref}{{[originallemma][0][]08DF}{[1][0][]source}{}{}{}}
\makeatother
\crefname{originallemma}{original lemma}{original lemmas}
\Crefname{originallemma}{Original lemma}{Original lemmas}
\makeatletter
\newlabel{lemma-induction-principle}{{08GP}{source}{}{sourcepointer-2}{}}
\newlabel{lemma-induction-principle@cref}{{[originallemma][0][]08GP}{[1][0][]source}{}{}{}}
\makeatother
\crefname{originallemma}{original lemma}{original lemmas}
\Crefname{originallemma}{Original lemma}{Original lemmas}
\makeatletter
\newlabel{lemma-lift-map-from-perfect-complex-with-support}{{09IV}{source}{}{sourcepointer-3}{}}
\newlabel{lemma-lift-map-from-perfect-complex-with-support@cref}{{[originallemma][0][]09IV}{[1][0][]source}{}{}{}}
\makeatother
\crefname{originallemma}{original lemma}{original lemmas}
\Crefname{originallemma}{Original lemma}{Original lemmas}
\makeatletter
\newlabel{lemma-pushforward-perfect}{{08HK}{source}{}{sourcepointer-4}{}}
\newlabel{lemma-pushforward-perfect@cref}{{[originallemma][0][]08HK}{[1][0][]source}{}{}{}}
\makeatother
\crefname{originallemma}{original lemma}{original lemmas}
\Crefname{originallemma}{Original lemma}{Original lemmas}
\makeatletter
\newlabel{lemma-pushforward-with-support-in-open}{{08GG}{source}{}{sourcepointer-5}{}}
\newlabel{lemma-pushforward-with-support-in-open@cref}{{[originallemma][0][]08GG}{[1][0][]source}{}{}{}}
\makeatother
\crefname{originallemma}{original lemma}{original lemmas}
\Crefname{originallemma}{Original lemma}{Original lemmas}
\makeatletter
\newlabel{perfect-lemma-generator-with-support}{{0A9A}{source}{}{sourcepointer-6}{}}
\newlabel{perfect-lemma-generator-with-support@cref}{{[originallemma][0][]0A9A}{[1][0][]source}{}{}{}}
\makeatother
\crefname{originallemma}{original lemma}{original lemmas}
\Crefname{originallemma}{Original lemma}{Original lemmas}
\makeatletter
\newlabel{perfect-lemma-orthogonal-koszul-complex}{{09IR}{source}{}{sourcepointer-7}{}}
\newlabel{perfect-lemma-orthogonal-koszul-complex@cref}{{[originallemma][0][]09IR}{[1][0][]source}{}{}{}}
\makeatother
\crefname{originallemma}{original lemma}{original lemmas}
\Crefname{originallemma}{Original lemma}{Original lemmas}
\makeatletter
\newlabel{remark-match-total-direct-images}{{08GH}{source}{}{sourcepointer-8}{}}
\newlabel{remark-match-total-direct-images@cref}{{[originallemma][0][]08GH}{[1][0][]source}{}{}{}}
\makeatother
\title{Stacks, Derived Categories of Spaces: generator-with-support}
\author{The Stacks Project Authors}
\date{stacks-project:spaces-perfect.tex:local-snapshot-20261012; GNU FDL 1.2 or later}
\begin{document}
\maketitle
\noindent\textit{Editorial note.} Statement span is lines 3996-4003 of sources/stacks-spaces-perfect/pinned-original-source.tex; proof span is lines 4005-4077, the complete author proof (71 body lines).\par\medskip
\noindent\textit{Editorial note.} Navigation stubs: 8 cited-but-not-included labels are supplied as non-load-bearing pointers with their pinned tags where the snapshot tag index has them.\par\medskip
\noindent\textit{Editorial note.} extra\_packages is empty on purpose: the retained preamble already loads xy with options and a second declaration clashes in the assembler wrapper.\par\medskip
\noindent\textit{Editorial note.} Text is retained verbatim, including the author's own typos and placeholders.\par\medskip
\noindent\textit{Editorial note.} Queue index 24 of queue\_spaceperfect.json (pb=3244, stmt 3996-4003, proof 4005-4077).\par\medskip
\noindent\textit{Editorial note.} stacks-spaces-perfect wave sp01 (2026-10-14): fresh adjudication accept per lane-s3/sp01\_judgments; pool per Chief 12:4x unfreeze (reserve-pool workflow, six-check clean); spaces-perfect.tex 126 with-proof rows, 70 b107-occupied labels avoided, 56 free rows judged 32 accept / 24 reject.\par\medskip
\noindent\textit{Allowed standard prerequisites (editorial).}\begin{itemize}
\item Explicit prerequisite boundary. The cited inputs are prerequisites; the argument in the retained span is not replaced by any cited result.
\end{itemize}\medskip
\ifdefined\hypertarget\hypertarget{sourcepointer-1}{}\fi
\noindent\textit{Source reference.} Original lemma 08DF: sources/stacks-spaces-perfect/pinned-original-source.tex, source lines 4077--4077. Cited by the retained statement or proof but not part of the collected spans; supplied only as a navigation pointer. The argument does not depend on its content.\par\medskip
\ifdefined\hypertarget\hypertarget{sourcepointer-2}{}\fi
\noindent\textit{Source reference.} Original lemma 08GP: sources/stacks-spaces-perfect/pinned-original-source.tex, source lines 1563--1577. Cited by the retained statement or proof but not part of the collected spans; supplied only as a navigation pointer. The argument does not depend on its content.\par\medskip
\ifdefined\hypertarget\hypertarget{sourcepointer-3}{}\fi
\noindent\textit{Source reference.} Original lemma 09IV: sources/stacks-spaces-perfect/pinned-original-source.tex, source lines 3814--3825. Cited by the retained statement or proof but not part of the collected spans; supplied only as a navigation pointer. The argument does not depend on its content.\par\medskip
\ifdefined\hypertarget\hypertarget{sourcepointer-4}{}\fi
\noindent\textit{Source reference.} Original lemma 08HK: sources/stacks-spaces-perfect/pinned-original-source.tex, source lines 3610--3616. Cited by the retained statement or proof but not part of the collected spans; supplied only as a navigation pointer. The argument does not depend on its content.\par\medskip
\ifdefined\hypertarget\hypertarget{sourcepointer-5}{}\fi
\noindent\textit{Source reference.} Original lemma 08GG: sources/stacks-spaces-perfect/pinned-original-source.tex, source lines 2135--2145. Cited by the retained statement or proof but not part of the collected spans; supplied only as a navigation pointer. The argument does not depend on its content.\par\medskip
\ifdefined\hypertarget\hypertarget{sourcepointer-6}{}\fi
\noindent\textit{Source reference.} Original lemma 0A9A: sources/stacks-spaces-perfect/pinned-original-source.tex, source lines 4077--4077. Cited by the retained statement or proof but not part of the collected spans; supplied only as a navigation pointer. The argument does not depend on its content.\par\medskip
\ifdefined\hypertarget\hypertarget{sourcepointer-7}{}\fi
\noindent\textit{Source reference.} Original lemma 09IR: sources/stacks-spaces-perfect/pinned-original-source.tex, source lines 4077--4077. Cited by the retained statement or proof but not part of the collected spans; supplied only as a navigation pointer. The argument does not depend on its content.\par\medskip
\ifdefined\hypertarget\hypertarget{sourcepointer-8}{}\fi
\noindent\textit{Source reference.} Original lemma 08GH: sources/stacks-spaces-perfect/pinned-original-source.tex, source lines 896--962. Cited by the retained statement or proof but not part of the collected spans; supplied only as a navigation pointer. The argument does not depend on its content.\par\medskip
\begin{lemma}
\label{lemma-generator-with-support}
Let $S$ be a scheme. Let $X$ be a quasi-compact and quasi-separated
algebraic space over $S$. Let $T \subset |X|$ be a
closed subset such that $|X| \setminus T$ is quasi-compact. With notation
as above, the category $D_{\QCoh, T}(\mathcal{O}_X)$ is generated by a
single perfect object.
\end{lemma}

\begin{proof}
We will prove this using the induction principle of
Lemma \ref{lemma-induction-principle}.
The property is true for representable quasi-compact
and quasi-separated objects of the site
$X_{spaces, \etale}$ by
Derived Categories of Schemes, Lemma
\ref{perfect-lemma-generator-with-support}.

\medskip\noindent
Assume that $(U \subset X, f : V \to X)$ is an elementary distinguished
square such that the lemma holds for $U$ and $V$ is affine. To finish
the proof we have to show that the result holds for $X$.
Let $P$ be a perfect object of $D(\mathcal{O}_U)$ supported on $T \cap U$
which is a generator for $D_{\QCoh, T \cap U}(\mathcal{O}_U)$. Using
Lemma \ref{lemma-lift-map-from-perfect-complex-with-support}
we may choose a perfect object $Q$ of $D(\mathcal{O}_X)$ supported on $T$
whose restriction to $U$ is a direct sum one of whose summands is $P$.
Write $V = \Spec(B)$. Let $Z = X \setminus U$. Then $f^{-1}Z$ is a closed
subset of $V$ such that $V \setminus f^{-1}Z$ is quasi-compact. As $X$
is quasi-separated, it follows that
$f^{-1}Z \cap f^{-1}T = f^{-1}(Z \cap T)$ is a closed
subset of $V$ such that $W = V \setminus f^{-1}(Z \cap T)$ is quasi-compact.
Thus we can choose $g_1, \ldots, g_s \in B$ such that
$f^{-1}(Z \cap T) = V(g_1, \ldots, g_r)$.
Let $K \in D(\mathcal{O}_V)$ be the perfect object corresponding to the
Koszul complex on $g_1, \ldots, g_s$ over $B$. Note that since $K$ is
supported on $f^{-1}(Z \cap T) \subset V$ closed, the pushforward
$K' = R(V \to X)_*K$ is a perfect object of $D(\mathcal{O}_X)$ whose
restriction to $V$ is $K$ (see Lemmas \ref{lemma-pushforward-perfect}
and \ref{lemma-pushforward-with-support-in-open}).
We claim that $Q \oplus K'$ is a generator for
$D_{\QCoh, T}(\mathcal{O}_X)$.

\medskip\noindent
Let $E$ be an object of $D_{\QCoh, T}(\mathcal{O}_X)$ such that
there are no nontrivial maps from any shift of $Q \oplus K'$ into $E$.
By Lemma \ref{lemma-pushforward-with-support-in-open}
we have $K' =  R(V \to X)_! K$ and hence
$$
\Hom_{D(\mathcal{O}_X)}(K'[n], E) = \Hom_{D(\mathcal{O}_V)}(K[n], E|_V)
$$
Thus by
Derived Categories of Schemes,
Lemma \ref{perfect-lemma-orthogonal-koszul-complex} we have
$E|_V = Rj_*E|_W$ where $j : W \to V$ is the inclusion. Picture
$$
\xymatrix{
W \ar[r]_j & V & Z \cap T \ar[l] \ar[d] \\
V \setminus f^{-1}Z \ar[u]^{j'} \ar[ru]_{j''} & & Z \ar[lu]
}
$$
Since $E$ is supported on $T$ we see that $E|_W$ is supported on
$f^{-1}T \cap W = f^{-1}T \cap (V \setminus f^{-1}Z)$
which is closed in $W$. We conclude that
$$
E|_V = Rj_*(E|_W) = Rj_*(Rj'_*(E|_{U \cap V})) = Rj''_*(E|_{U \cap V})
$$
Here the second equality is part (1) of
Cohomology, Lemma \ref{cohomology-lemma-pushforward-restriction}
which applies because $V$ is a scheme and $E$ has quasi-coherent
cohomology sheaves hence pushforward along the quasi-compact
open immersion $j'$ agrees with pushforward on the underlying schemes, see
Remark \ref{remark-match-total-direct-images}.
This implies that $E = R(U \to X)_*E|_U$ (small detail omitted). If
this is the case then
$$
\Hom_{D(\mathcal{O}_X)}(Q[n], E) = \Hom_{D(\mathcal{O}_U)}(Q|_U[n], E|_U)
$$
which contains $\Hom_{D(\mathcal{O}_U)}(P[n], E|_U)$ as a direct summand.
Thus by our choice of $P$ the vanishing of these groups implies that $E|_U$
is zero. Whence $E$ is zero.
\end{proof}

\end{document}
